% 日本語草稿用の preamble（XeLaTeX）。投稿先の class や style に合わせて書き換えてよい。
\usepackage{fontspec}
\IfFontExistsTF{Noto Serif CJK JP}{%
  \setmainfont{Noto Serif CJK JP}\setsansfont{Noto Sans CJK JP}%
}{\IfFontExistsTF{IPAexMincho}{%
  \setmainfont{IPAexMincho}\setsansfont{IPAexGothic}%
}{\IfFontExistsTF{Hiragino Mincho ProN}{%
  \setmainfont{Hiragino Mincho ProN}\setsansfont{Hiragino Sans}%
}{}}}
\XeTeXlinebreaklocale "ja"
\XeTeXlinebreakskip=0pt plus 1pt
\usepackage[margin=25mm]{geometry}
\usepackage{amsmath,amssymb}
\usepackage{graphicx}
\usepackage{booktabs}
\usepackage[authoryear,round]{natbib}
\usepackage{xcolor}
\usepackage[unicode,colorlinks=true,linkcolor=blue,citecolor=blue,urlcolor=blue]{hyperref}

% 図は figures/<name>/<name>.pdf にある: \includegraphics{<name>/<name>}
\graphicspath{{../figures/}}

\renewcommand{\abstractname}{要旨}
\renewcommand{\refname}{参考文献}
\renewcommand{\figurename}{図}
\renewcommand{\tablename}{表}

% PDF に見える形で残す著者向けのメモ（要出典、確認待ちなど）。
\newcommand{\note}[1]{\textcolor{blue}{[Note: #1]}}

% まだ得ていない結果の予想。数値にも文にも使い、結果が出たら外す。
\newcommand{\pred}[1]{\textcolor{orange!85!black}{#1}}
